\subsection{Определения углов.}
В данном разделе мы рассмотрим два определения угла облака. Для того, чтобы
дать эти определения сначала необходимо определить угол между пространствами.
\\
В дальнейшем будем считать, что у облаков нетривиальная стационарная группа и в
обозначении облака $ [M] $, $ M $ является его центром.
\begin{definition}
  \emph{Углом} между пространствами $ X_1, X_2 \in [M] $ такими, что $ |X_1
  M| = r_1, | X_2 M | =r_2, |X_1 X_2| = d$, где $ r_1,
  r_2 \neq 0 $ называется величина $ \arccos \left(\frac{r_1^2 +
  r_{2}^2 - d^2}{2r_1r_2} \right)$.\\ Будем обозначать его $
  \varphi(X_1, X_2) $.
\end{definition}
\begin{remark}
  Такое определение естественным образом вытекает из евклидовой теоремы
  косинусов, а именно $ c^{2} = a^2 + b^2 - 2ab\cos (a,b) $.
\end{remark}
Прежде, чем формулировать свойство угла, нам понадобится следующая
вспомогательная лемма.
\begin{lemma}
  \label{lem:lambdacloud} Для любого облака \( [M] \) и \( \lambda
  \in\mathbb{R}^+ \) выполняются следующие свойства{\textup{:}}
  \begin{enumerate}
    \item \( \St\big(\lambda [M]\big) = \St\big([M]\big)\),
      \label{lem:lambdacloud:1}
    \item \( \lambda [M] = [\lambda M] \).\label{lem:lambdacloud:2}
  \end{enumerate}
\end{lemma}
\begin{proof}
  (\ref{lem:lambdacloud:1}) Пусть \( \mu \in \St\big([M]\big) \). Тогда \( \mu
  \cdot \lambda [M] = \lambda \cdot \mu [M]= \lambda [M] \), т. е. \( \mu \in
  \St\big(\lambda [M]\big) \). Отсюда \( \St\big([M]\big) \subseteq
  \St\big(\lambda [M]\big) \). Обратное включение получается если рассмотреть
  облако \( [M] \) как \(\frac{1}{\lambda}\cdot \lambda [M] \).

  (\ref{lem:lambdacloud:2}) Так как стационарные группы совпадают, для любого
  \( \mu \in  \St\big(\lambda [M])\) выполняется \( \mu \cdot \lambda M =
  \lambda \cdot \mu M = \lambda M \). Значит, \( \lambda M \) --- центр
  облака \( \lambda [M] \).
\end{proof}
У угла между пространствами есть следующее свойство.
\begin{lemma}
  \label{lemmaAngleBetweenSpaces} Для любых \( X_1, X_2 \in [M], \lambda
  \in\mathbb{R}^+ \) выполняется \( \varphi (X_1, X_2) = \varphi (\lambda
      X_1,
  \lambda X_2) \).
\end{lemma}
\begin{proof}
  По лемме \ref{lem:lambdacloud} пространство \( \lambda M \) будет
  центром \( \lambda
    [M]
  \). Как в определении угла, положим \[ |X_1 M| = r_1, | X_2, M | =r_2,
    |X_1
  X_2| = d.\] По определению угла получаем следующую цепочку
  равенств: \[ \varphi (\lambda X_1, \lambda X_2) = \frac{\lambda
    ^{2} r_1^2 + \lambda ^2 r_2^2 - \lambda^2 d^2}{2 \lambda  r_1
    \lambda r_2} =
    \varphi (X_1, X_2).
  \]
  На этом доказательство закончено.
\end{proof}
Рассмотрим теперь два интересующих нас определения угла облака.
\begin{definition}
  \emph{Углом} облака $ [M] $ называется величина
  \[
    \varphi \big([M]\big)= \sup \big\{\varphi (X_1, X_2) \colon | X_1,M
    |, |X_2,M| \neq 0\big\}.
  \]
\end{definition}
\begin{definition}
  \emph{Равнобедренным углом} облака $ [M] $ называется величина
  \[
    \varphi_e \big([M]\big) = \sup \big\{\varphi (X_1, X_2) \colon |X_1,M |
    = |X_2,M| \neq 0\big\}.
  \]
\end{definition}
Для этих определений получаем следствие из леммы \ref{lemmaAngleBetweenSpaces}.
\begin{lemma}
  Для любого облака \( [M] \) и \( \lambda \in\mathbb{R}^+ \)
  выполняется\textup{:}\label{lem:angles}
  \begin{enumerate}
    \item \(\varphi \big(\lambda [M]\big) = \varphi
      \big([M]\big)\),\label{lem:angles:1}
    \item \(\varphi_e \big(\lambda [M]\big) = \varphi_e
      \big([M]\big)\),\label{lem:angles:2}
    \item \( \varphi_e \big([M]\big) \le \varphi \big([M]\big)
      \),\label{lem:angles:3}
    \item \(0 \le \varphi ([M]) \le \pi\),\label{lem:angles:4}
    \item \(0 \le \varphi_e ([M]) \le \pi\).\label{lem:angles:5}
  \end{enumerate}
\end{lemma}

Приведем примеры углов облаков.
\begin{lemma}
  Для облаков \( [\Delta _{1}], [\mathbb{R}] \) выполняются следующие
  равенства: \label{lem:angleexam}
  \begin{enumerate}
    \item \( \varphi \big([\Delta _{1}]\big) = \frac{\pi }{2}
      \),\label{lem:angleexam:1}
    \item \( \varphi_e \big([\Delta _{1}]\big) = \frac{\pi }{3}
      \),\label{lem:angleexam:2}
    \item \( \varphi \big([\mathbb{R}]\big) = \varphi_e
      \big([\mathbb{R}]\big) = \pi  \).\label{lem:angleexam:3}
  \end{enumerate}
\end{lemma}
\begin{proof}
  (\ref{lem:angleexam:1}) Из ультраметрического неравенства следует, что \(
  \varphi \big([\Delta _1]\big)\le \frac{\pi }{2} \). Для доказательства
  обратного неравенства найдем угол между двухточечным симплексом \( \Delta
  _{2} = \{x_1, x_2\} \) и трехточечным симплексом \(
  \lambda \Delta _3  = \{y_1,y_2, y_3\}\). Воспользуемся формулой, доказанной в
  \cite{ivanov2016} для нахождения расстояния Громова--Хаусдорфа
  между двумя трехточечными пространствами:
  \[
    d_{GH}(X,Y) = \frac{1}{2} \max\big\{ |a_1 - a_2|,|b_1 - b_2|,|c_1
    - c_2|\big\}.
  \]
  Здесь \( a_1 \le b_1 \le c_1 \) --- расстояния между точками в
  пространстве \( X \), \( a_2 \le b_2 \le c_2 \) --- расстояния между
  точками в пространстве \( Y \). В нашем случае это \( 0,1,1 \) и \(
  \lambda , \lambda , \lambda . \) Соответственно,
  \[
    d_{GH}(\Delta_2, \lambda \Delta _{3}) =\frac{1}{2}
    \max\big\{|\lambda |, | \lambda  -1|, |\lambda -1|\big\},
  \]
  что
  при \( \lambda \ge \frac 1 2 \) равно \(\frac \lambda 2 \). Итак, \(d
    _{GH}(\Delta _2, \lambda \Delta _3) =
  \frac{\lambda }{2} \). Далее, учитывая, что \( |\Delta _1
    \Delta _2|=\frac{1}{2}, |\Delta _1 \lambda \Delta
  _3|=\frac{\lambda }{2} \), найдем угол между \( \Delta _2,
  \lambda \Delta _3 \):
  \[ \cos \big(\varphi (\Delta _2 , \lambda \Delta _3)\big)
    =\frac{\frac{\lambda ^2}{4} + \frac{1}{4} - \frac{\lambda
    ^2}{4}}{\frac{\lambda }{2}} = \frac{1}{2 \lambda }
    \rightarrow 0, \lambda \rightarrow \infty.
  \]
  Отсюда, \( \varphi \big([\Delta _1]\big) \ge \sup \varphi (\Delta _2,
  \lambda \Delta _3) = \frac \pi 2. \)

  (\ref{lem:angleexam:2}) Из ультраметрического неравенства следует, что \(
  \varphi_e \big([\Delta _1]\big)\le \frac{\pi }{3} \). Для доказательства
  обратного неравенства достаточно найти пространства одинакового диаметра,
  расстояние между которыми равно их диаметрам. Этими пространствами
  являются, например, \( \Delta _2, \Delta _3 \).

  (\ref{lem:angleexam:3}) Достаточно привести пример пространств \( X,Y\in
  [\mathbb{R}] \), таких, что \( |X \mathbb{R}| = |Y \mathbb{R}| = \frac 1 2
  |XY|. \) В \cite{mikhailovUnreleased} было доказано, что такими
  пространствами являются \( \mathbb{Z}, \mathbb{R}\times
  [0,1] \).
  %Вставить ссылку на Ивана с названием.

\end{proof}
\subsection{Расстояние между облаками с разными углами.}
\begin{theorem}
  Пусть облака \( [M], [N] \) имеют нетривиальное пересечение
  стационарных групп и их углы \( \varphi \big([M]\big), \varphi
  \big([N]\big) \) различны. Тогда \( d_{GH} \big([M], [N]\big) =
  \infty \).
\end{theorem}
\begin{proof}
  Без ограничения общности будем считать, что \( \varphi \big([M]\big)>
  \varphi \big([N]\big) \). Для доказательства теоремы покажем, что не
  существует соответствия \( R\in \mathcal{R}\big([M],[N]\big) \) с
  конечным искажением.

  Предположим противное, т. е. существует соответствие \( R \) с
  конечным искажением \( \dis R = \varepsilon < \infty \) как
  отображение из облака \( [M] \) в облако \( [N] \). Построим функцию
  \( f \colon [M] \rightarrow [N] \), \( f(X) \) -- произвольное
  пространство из \(
    R(X)
  \). По определению угла, в облаке \( [M] \) найдутся пространства \(
  X_1, X_2 \), для которых выполнено
  \[
    \varphi \big([N]\big) < \varphi (X_1, X_2) \le \varphi
  \big([M]\big). \]
  Положим \( |X_1 M| = r_1^X,  |X_2 M| = r_2^X, | X_1 X_2 | = d^X\).
  Также зададим функцию
  \[ d^N \colon \mathbb{R}^+ \times \mathbb{R}^+ \rightarrow
    \mathbb{R}^+, d^N(r_1, r_2) =
    \sqrt{r_1^2 + r_2^2 - 2r_1r_2\cos\Big(\varphi \big([N]\big)\Big)}.
  \]
  Тогда выполняется следующее неравенство:
  \begin{equation}
    d^N(r_1^X, r_2^X) < d^X.
    \label{ineqdx}
  \end{equation}
  По определению угла для любых пространств \( Y_1, Y_2 \in [N] \)
  выполняется
  \begin{equation}
    | Y_1 Y_2 | \le d^N\big(| Y_1N |, | Y_2N |\big). \label{ineqAngle}
  \end{equation}
  Итак, наша задача --- показать, что найдутся пространства в \( [N]
  \), для которых неравенство (\ref{ineqAngle}) не выполняется.

  Нетривиальность пересечения стационарных групп означает, что найдется
  подгруппа \[
    \{q^k \colon k\in \mathbb{Z}, q> 1\} \in \St\big([M]\big) \cap \St
    \big([N]\big).
  \]
  Положим \( |N R(M)| = l \). Будем рассматривать пространства \(
  R(q^k X_1), R(q^k X_2), k \in \mathbb{N} \). Для них выполнены
  следующие неравенства:
  \begin{equation}
    \big| R(q^k X_1), R(q^k X_2) \big| \ge q^k d^X - \varepsilon,
    \label{ineqDistx1x2}
  \end{equation}
  \[
    q^k r_1^X - \varepsilon  \le\big | R(M), R(q^k X_1) \big| \le q^k
    r_1^X + \varepsilon ,
  \]
  \[
    q^k r_2^X - \varepsilon  \le\big | R(M), R(q^k X_2) \big| \le q^k
    r_2^X + \varepsilon .
  \]
  Из последних двух неравенств по неравенству треугольника получаем
  следующее:
  \begin{equation}
    q^k r_1^X - l - \varepsilon \le \big|N, R(q^k X_1) \big| \le
    q^k r_1^X + l + \varepsilon,\label{ineqDistx1N}
  \end{equation}
  \begin{equation}
    q^k r_2^X - l - \varepsilon \le \big|N, R(q^k X_2) \big| \le q^k
    r_2^X + l + \varepsilon.\label{ineqDistx2N}
  \end{equation}
  Поделим неравенства (\ref{ineqDistx1x2}), (\ref{ineqDistx1N}),
  (\ref{ineqDistx2N}) на \( q^k \):
  \[ \big| q^{-k}R(q^k X_1), q^{-k}R(q^k X_2) \big| \ge d^X -
  \frac{\varepsilon }{q^k}, \]
  \[
    r_1^X - \frac{l + \varepsilon}{q^k} \le \big|N, q^{-k}R(q^k X_1)
    \big| \le
    r_1^X + \frac{l + \varepsilon}{q^k},
  \]
  \[
    r_2^X - \frac{l + \varepsilon}{q^k} \le \big|N, q^{-k}R(q^k X_2)
    \big| \le
    r_2^X + \frac{l + \varepsilon}{q^k}.
  \]
  Функция \( d^N \) непрерывная, значит
  \[
    \lim_{k \rightarrow \infty
    }d^N\Big(\big|N, q^{-k}R(q^k X_1)
      \big|, \big|N, q^{-k}R(q^k X_2)
    \big|\Big) = d^N(r_1^X, r_2^X).
  \]
  Из неравенства (\ref{ineqdx}) следует, что найдутся \( \delta
    _1, \delta _2
  > 0 \) такие, что \( d^N(r_1^X,r_2^X)<d^N(r_1^X, r_2^X)+\delta
  _1 < d^X- \delta _2 < d^X \). Выберем такое \( k_1 \), чтобы
  для всех \( k > k_1 \) выполнялось
  \[
    d^N\Big(\big|N, q^{-k}R(q^k X_1)
      \big|, \big|N, q^{-k}R(q^k X_2)
    \big|\Big) \in \big(d^N(r_1^X, r_2^X)- \delta _1,
  d^N(r_1^X, r_2^X)+ \delta _1\big).  \]
  Также выберем \( k_2 \), такое что \( \big| q^{-k}R(q^k X_1),
  q^{-k}R(q^k X_2) \big| > d^X - \delta _2 \).
  Теперь, если взять \( k_0 = \max\{k_1,k_2\} \), то
  \[ \big| q^{-k_0}R(q^{k_0} X_1), q^{-k_0}R(q^{k_0} X_2) \big|
    > d^N\Big(\big|N, q^{-k_0}R(q^k_0 X_1) \big|, \big|N,
  q^{-k_0}R(q^{k_0} X_2) \big|\Big) \]
  Пространства \( q^{-k_0}R(q^{k_0}X_1), q^{-k_0}R(q^{k_0}X_2)\)
  ---
  искомые, для которых не выполняется неравенство (\ref{ineqAngle}).
  Противоречие.
\end{proof}
